\cleardoublepage
\chapter{ANALISIS MASALAH}
\label{chap:analisis-masalah}

Bab ini membahas analisis kondisi operasional TransJogja saat ini, kebutuhan
yang harus dipenuhi simulasi untuk menjawab rumusan masalah pada Bab I, serta
analisis alternatif solusi yang dipertimbangkan sebelum menentukan pendekatan
akhir yang digunakan pada Bab IV dan Bab V.

\section{Analisis Kondisi Saat Ini}

Gambar~\ref{fig:kondisi-saat-ini} menunjukkan model konseptual kondisi
operasional TransJogja saat ini beserta komponen dan interaksi
antarkomponennya.

\begin{figure}[h]
	\centering
	\begin{tikzpicture}[
		node distance=1.3cm and 1.7cm,
		box/.style={draw, rounded corners, minimum width=3.2cm, minimum height=1cm, align=center, font=\small},
		busbox/.style={box, fill=blue!8},
		infra/.style={box, fill=orange!12},
		halte/.style={draw, circle, minimum size=0.5cm, fill=gray!15, font=\scriptsize},
		arrow/.style={-Stealth, thick},
		lbl/.style={font=\scriptsize\itshape}
		]
		\node[halte] (h1) {J};
		\node[halte, right=1.6cm of h1] (h2) {M};
		\node[halte, right=1.6cm of h2] (h3) {J};
		\draw[arrow] (h1) -- (h2) node[midway, above, lbl] {8,22 km};
		\draw[arrow] (h2) -- (h3) node[midway, above, lbl] {8,22 km};
		\node[lbl, below=0.05cm of h1] {Term. Jombor};
		\node[lbl, below=0.05cm of h2] {Malioboro};
		\node[lbl, below=0.05cm of h3] {Term. Jombor};
		\node[busbox, below=1.9cm of h1, xshift=0.3cm] (bus) {2 Bus Listrik\\Rute L-1};
		\node[infra, right=2.6cm of bus] (spklu) {SPKLU\\Adisucipto\\(1 lokasi)};
		\draw[arrow] (bus.east) -- node[above, lbl]{mengisi daya} (spklu.west);
		\draw[arrow] (spklu.north) to[bend left=35] node[above, lbl, pos=0.5]{kembali beroperasi} (h2.south);
		\node[box, fill=gray!5, below=1.6cm of bus, xshift=1.3cm, minimum width=6.8cm] (dishub) {Dishub DIY\\(menetapkan target headway per koridor)};
		\draw[arrow] (dishub.north) -- (bus.south);
		\begin{scope}[on background layer]
			\node[draw, dashed, fit=(h1)(h2)(h3)(bus)(spklu)(dishub), inner sep=0.4cm] {};
		\end{scope}
	\end{tikzpicture}
	\caption{Model konseptual kondisi operasional TransJogja saat ini}
	\label{fig:kondisi-saat-ini}
\end{figure}

TransJogja saat ini mengoperasikan 2 unit bus listrik pada Rute L-1
(Terminal Jombor, Malioboro, Terminal Jombor), dengan satu SPKLU di kawasan
Adisucipto sebagai satu-satunya titik pengisian daya. Rute L-1 adalah
satu-satunya data lapangan yang tersedia untuk kondisi bus listrik
TransJogja: bus mengisi daya secara \textit{opportunity charging} (mengisi
daya di sela rit, bukan hanya di depot), dengan rata-rata 2 sesi pengisian
per hari dan total waktu tidak beroperasi (\textit{downtime}) sekitar 120
menit per sesi, mencakup waktu tempuh menuju SPKLU maupun waktu pengisian
itu sendiri.

Kondisi ini berbeda dari rencana Dishub DIY untuk mengelektrifikasi seluruh
20 koridor TransJogja (kira-kira 124 sampai 128 bus). Pada skala 2 bus dan
1 rute, penempatan SPKLU tunggal di Adisucipto tidak menimbulkan persoalan
berarti karena hanya 1 rute yang bergantung padanya. Pada skala penuh,
setiap koridor punya jarak dan pola permintaan pengisian daya yang
berbeda-beda, sehingga jumlah dan lokasi SPKLU yang memadai untuk 1 koridor
belum tentu memadai untuk koridor lain. Dokumen perencanaan Dishub DIY saat
ini menetapkan target headway per koridor, tetapi belum menyertakan analisis
tentang bagaimana target itu tetap tercapai setelah elektrifikasi penuh.
Kesenjangan inilah yang dianalisis lebih lanjut pada subbab berikut.

\section{Analisis Kebutuhan}

\subsection{Identifikasi Masalah Pengguna}

Terdapat dua kelompok pemangku kepentingan utama yang kebutuhannya
memengaruhi rancangan simulasi ini:
\begin{enumerate}
	\item \textbf{Dishub DIY/operator (PT AMI)} sebagai calon pengambil
	keputusan investasi infrastruktur. Kebutuhan mereka adalah pembanding
	kuantitatif antar-skenario lokasi, kapasitas, dan mekanisme pengisian
	daya, disertai dampaknya terhadap headway operasional. Ini penting karena
	rekomendasi kualitatif saja sulit dipertanggungjawabkan sebelum investasi
	dilakukan.
	\item \textbf{Penulis (sebagai peneliti/pengembang)} yang memerlukan model
	dan mesin simulasi yang dapat divalidasi terhadap data lapangan yang
	tersedia (baseline Rute L-1), agar hasil simulasi untuk 20 koridor
	fully-electric dapat dipertanggungjawabkan meski belum ada data lapangan
	langsung untuk skala tersebut.
\end{enumerate}

Dari kedua kebutuhan ini, masalah pengguna dapat dirumuskan sebagai:
bagaimana menyediakan alat bantu simulasi yang (a) cukup valid untuk
dipercaya sebagai dasar pengambilan keputusan, dan (b) cukup jelas hasilnya
untuk dibaca dan dibandingkan oleh pihak yang tidak terlibat langsung dalam
pengembangan model.

\subsection{Kebutuhan Fungsional}

Berdasarkan identifikasi masalah pengguna di atas, kebutuhan fungsional
sistem simulasi adalah sebagai berikut:
\begin{enumerate}
	\item Sistem dapat memodelkan drain baterai non-linear untuk tiap bus
	berdasarkan jarak antar-halte dan kecepatan tempuh per segmen rute.
	\item Sistem dapat menjalankan simulasi operasi seluruh 20 koridor secara
	simultan, termasuk penjadwalan rit tiap bus sepanjang jam operasional.
	\item Sistem dapat memodelkan antrean pengisian daya di tiap SPKLU sebagai
	\textit{resource} dengan kapasitas terbatas (jumlah SPKLU per lokasi).
	\item Sistem dapat menjalankan skenario lokasi SPKLU yang berbeda-beda
	(kondisi saat ini, rencana yang sudah ada, dan kandidat lokasi baru)
	sebagai konfigurasi \textit{input} yang dapat dibandingkan.
	\item Sistem dapat menjalankan skenario mekanisme pengisian daya yang
	berbeda-beda (\textit{opportunity charging}, \textit{depot charging}, dan
	kombinasi keduanya), termasuk opsi penambahan bus per koridor.
	\item Sistem dapat menghitung dan membandingkan headway aktual hasil
	simulasi terhadap target headway pada dokumen perencanaan Dishub DIY,
	per koridor dan per skenario.
	\item Sistem dapat menjalankan repetisi simulasi (Monte Carlo) dengan
	variasi waktu tempuh acak antar-hari, untuk menangkap variasi kondisi
	operasional harian yang tidak dapat ditangkap satu kali simulasi
	deterministik.
	\item Sistem menyediakan dasbor visual yang menampilkan peta skematik rute
	dan animasi pergerakan bus per skenario, agar hasil dapat dibaca tanpa
	perlu memeriksa data mentah.
\end{enumerate}

\subsection{Kebutuhan Nonfungsional}

\begin{enumerate}
	\item \textbf{Validitas:} model drain baterai dan mesin simulasi harus
	dapat divalidasi secara kualitatif terhadap data operasional riil
	Rute L-1 sebagai baseline, sebelum digunakan untuk skenario 20 koridor
	yang belum memiliki data lapangan.
	\item \textbf{Skalabilitas komputasi:} mesin simulasi harus dapat
	dijalankan secara \textit{headless} (tanpa antarmuka) dalam mode batch
	untuk puluhan skenario dan ratusan repetisi Monte Carlo per skenario,
	tanpa memerlukan waktu komputasi yang tidak wajar untuk pengerjaan tugas
	akhir.
	\item \textbf{Keterpisahan (\textit{separation of concerns}):} logika
	simulasi harus terpisah dari lapisan visualisasi, sehingga dasbor cukup
	membaca hasil simulasi yang sudah tersimpan tanpa perlu menjalankan ulang
	logika simulasi di sisi klien.
	\item \textbf{Reprodusibilitas:} setiap repetisi Monte Carlo harus dapat
	diulang dengan hasil yang identik (melalui \textit{random seed} yang
	konsisten dan terdokumentasi), agar hasil simulasi dapat diperiksa ulang.
	\item \textbf{Keterbacaan hasil:} dasbor visual harus dapat dikenali sebagai
	representasi peta rute TransJogja oleh pembaca yang familier dengan rute
	aslinya, agar hasil simulasi mudah dikomunikasikan ke pemangku
	kepentingan non-teknis.
\end{enumerate}

\section{Analisis Pemilihan Solusi}

\subsection{Alternatif Solusi}


%TODO-CITE: subbab ini dapat diperkuat dengan sitasi literatur mengenai
% model konsumsi energi kendaraan listrik dan/atau metode Discrete Event
% Simulation, untuk menunjukkan bahwa alternatif yang dipertimbangkan
% memang dikenal dalam studi terkait, bukan hanya pilihan ad-hoc.

Beberapa pertanyaan desain memiliki lebih dari satu alternatif solusi yang
dipertimbangkan sebelum solusi akhir ditentukan:

\begin{enumerate}
	\item \textbf{Model drain baterai.} Alternatif pertama adalah menyusun
	formula heuristik sendiri (koefisien \textit{k1}, \textit{k2}) yang
	dikalibrasi langsung terhadap baseline Rute L-1. Alternatif kedua adalah
	mengadopsi model konsumsi energi kendaraan listrik dari literatur yang
	sudah divalidasi pada kendaraan sejenis.
	\item \textbf{Arsitektur mesin simulasi.} Alternatif pertama adalah simulasi
	berbasis langkah waktu tetap (\textit{tick-based}), yaitu kondisi ketika status seluruh
	bus diperbarui pada interval waktu yang sama. Alternatif kedua adalah
	simulasi berbasis peristiwa diskrit (\textit{Discrete Event Simulation}),
	yaitu pembaruan status yang hanya dilakukan saat terjadi peristiwa (mis. bus tiba di
	halte, bus mulai/selesai mengisi daya).
	\item \textbf{Penentuan lokasi kandidat SPKLU.} Alternatif pertama adalah
	algoritma optimasi \textit{greenfield} yang secara otomatis mencari lokasi
	optimal dari data mentah (mis. pendekatan metaheuristik atau
	\textit{reinforcement learning}). Alternatif kedua adalah membandingkan
	sejumlah skenario lokasi yang ditentukan sebagai \textit{input}
	(\textit{scenario comparison}/\textit{grid search}), termasuk kandidat
	yang disusun dari analisis data permintaan pengisian daya hasil simulasi
	itu sendiri.
	\item \textbf{Mekanisme pengisian daya yang disimulasikan.} Alternatif yang
	dipertimbangkan mencakup \textit{opportunity charging} penuh, \textit{depot
		charging} penuh, kombinasi keduanya (\textit{mix}), dan mekanisme
	pengisian daya dinamis yang menyesuaikan target SoC berdasarkan kondisi
	antrean atau waktu.
	\item \textbf{Teknologi dasbor visual.} Alternatif pertama adalah dasbor
	berbasis \textit{framework} \textit{frontend} modern (React) dengan
	\textit{build toolchain} terpisah. Alternatif kedua adalah dasbor berbasis
	HTML/JavaScript murni yang disajikan langsung oleh \textit{backend} tanpa
	proses \textit{build} tambahan.
\end{enumerate}

\subsection{Analisis Penentuan Solusi}

\textbf{Model drain baterai.} Formula heuristik \textit{k1}/\textit{k2} yang
dikalibrasi langsung terhadap baseline Rute L-1 digunakan pada tahap awal
pengerjaan karena kesesuaiannya yang baik dengan data lapangan yang tersedia
(2 sesi pengisian per hari, sesuai kondisi riil). Namun, pendekatan ini
memiliki kelemahan mendasar: parameter \textit{k1}/\textit{k2} dikalibrasi
langsung dari satu baseline (Rute L-1) tanpa dasar fisik yang independen,
sehingga validitasnya untuk kondisi 20 koridor lain (jarak, kecepatan, dan
profil rute yang berbeda-beda) tidak terjamin. Solusi akhir yang dipilih
adalah model konsumsi energi non-linear dari \textcite{mamarikas2025}
(Persamaan 1 pada Bab IV), yang memiliki dasar fisik dan telah dikalibrasi
pada kendaraan listrik kelas berat pada studi aslinya. Konsekuensinya, model
ini menghasilkan konsumsi energi yang lebih tinggi dari kondisi riil bus
TransJogja (bus kelas menengah), sehingga baseline Rute L-1 bergeser dari
2 menjadi 3 sesi pengisian per hari pada hasil simulasi. Deviasi ini
diterima secara sadar sebagai batasan tugas akhir (lihat Bab I, Batasan
Masalah), karena mempertahankan integritas model literatur dinilai lebih
dapat dipertanggungjawabkan dibanding menyesuaikan hasil dengan faktor
skala buatan yang tidak berdasar data primer.

\textbf{Arsitektur mesin simulasi.} Pendekatan \textit{tick-based} pada
tahap awal pengembangan terbukti tidak efisien dan rentan terhadap
kesalahan logika (di antaranya bus yang "berteleportasi" antar-halte dan
penghitungan waktu kembali yang tidak akurat) ketika diterapkan pada 20
koridor secara simultan dengan jumlah bus yang besar. Solusi akhir yang
dipilih adalah arsitektur \textit{Discrete Event Simulation} berbasis
antrean prioritas (\textit{heapq}), yang hanya memproses peristiwa yang
benar-benar terjadi. Pendekatan ini lebih efisien secara komputasi untuk
skenario dengan banyak entitas (bus, SPKLU) yang aktivitasnya tersebar
tidak merata sepanjang waktu, dan menghilangkan kelas kesalahan yang
muncul pada pendekatan \textit{tick-based} sebelumnya.

\textbf{Penentuan lokasi kandidat SPKLU.} Algoritma optimasi otomatis
tidak dipilih karena berada di luar batasan masalah yang disepakati sejak
awal (lihat Bab I, Batasan Masalah). Tugas akhir ini berfokus pada
perbandingan skenario, bukan pengembangan algoritma optimasi. Solusi yang
dipilih adalah perbandingan skenario lokasi, dengan kandidat lokasi disusun
dari dua sumber: kondisi saat ini/rencana yang sudah ada, dan analisis data
permintaan pengisian daya yang diekstraksi dari hasil simulasi tahap
sebelumnya (\textit{data-driven}). Pendekatan ini dipilih karena tetap
memberi dasar kuantitatif untuk pemilihan lokasi tanpa memperluas cakupan
tugas akhir ke pengembangan algoritma optimasi baru.

\textbf{Mekanisme pengisian daya yang disimulasikan.} Keempat mekanisme
(\textit{opportunity}, \textit{depot}, \textit{mix}, dan dinamis)
diimplementasikan dan diuji, karena masing-masing menjawab pertanyaan
evaluasi yang berbeda: \textit{depot charging} penuh diuji untuk
mendemonstrasikan secara empiris mengapa \textit{opportunity charging}
diperlukan (skenario \textit{depot charging} penuh menghasilkan kegagalan
pada seluruh koridor), sementara mekanisme dinamis diuji untuk menjawab
apakah algoritma penjadwalan pengisian daya yang lebih adaptif memberi
manfaat tambahan di atas infrastruktur dasar yang direkomendasikan. Hasil
pengujian tiap mekanisme dijelaskan pada Bab VI.

\textbf{Teknologi dasbor visual.} Dasbor berbasis React sempat
dipertimbangkan pada tahap perancangan awal, namun tidak dipilih karena
menambah kerumitan \textit{build toolchain} yang tidak sebanding dengan
kebutuhan dasbor ini (menampilkan hasil simulasi yang sudah tersimpan,
bukan aplikasi interaktif kompleks). Solusi akhir yang dipilih adalah
\textit{backend} Flask (Python) yang menyajikan data hasil simulasi (JSON),
dikonsumsi oleh \textit{frontend} HTML/JavaScript murni dengan Chart.js
untuk visualisasi grafik. Pendekatan ini konsisten dengan kebutuhan
nonfungsional keterpisahan (lihat Analisis Kebutuhan) dan mempercepat
iterasi pengembangan karena tidak memerlukan proses \textit{build} terpisah.